% Generated by paper/latex_build.py from paper/*.md. Do not edit by hand.

This study informs how language models are compared and chosen. Its main practical message,
that the choice of model mattered more than the reasoning setting on this task, could be
over-read in two ways. First, it holds for one open-ended design task, six families and one
26-hour window; it is not evidence that reasoning modes are unhelpful in general, and on
some families the reasoning setting changed scores a great deal. Second, the per-family
scores are not a product ranking. They reflect one prompt, app settings that differed
between vendors (including live web search in three apps) and models that the apps can
change without notice. We report them to explain the variance decomposition, not to
recommend a vendor.

Two findings bear on evaluation practice beyond this paper. Both LLM raters scored their
own vendor's answers higher than the other raters did, by about as much as the average
effect of reasoning mode, so a study that uses a single LLM judge from a vendor it is
comparing can mistake self-preference for an effect. And two of six models named a
different vendor's model when asked to identify themselves, so a model's self-report is not
reliable evidence of which model answered.

Reasoning modes also use more computation per answer: in our API checks the reasoning
setting used between 1,367 and 23,805 reasoning tokens per call, where the fast setting used
0 to 131. Where reasoning does not improve the answer, the extra cost and energy buy little,
which is one more reason to report the mode.

The study collected no personal data. The 60 answers are model outputs to a fixed prompt
about computer hardware. The human rater is not named, at their request.
